\subsection{模型假设}

\begin{enumerate}
    \item \label{as:road} 实际道路距离与球面直线距离之间满足固定的区域比例关系。将全国城市按东部、中部、西部分别设定绕行系数 $\beta_E = 1.25$、$\beta_M = 1.40$、$\beta_W = 1.60$，同一区域内所有城市共享相同绕行系数。
    \item \label{as:speed} 车辆在高速公路上以平均速度 $v_h = 70\,\mathrm{km/h}$ 匀速行驶（问题一、二），在含城市道路的混合路况下以 $v_c = 60\,\mathrm{km/h}$ 匀速行驶（问题三），忽略交通拥堵、天气等随机因素。
    \item \label{as:load} 每次装卸货作业固定耗时 $t_{\mathrm{lu}} = 2\,\mathrm{h}$，装卸货工作仅能在每日9:00至17:00之间进行，单日最多可完成4次装卸操作。
    \item \label{as:empty} 车辆回程空载时油耗为满载油耗的75\%（空载系数 $\alpha = 0.75$），路桥费与去程相同，人工等固定成本不受负载影响。
    \item \label{as:week} 问题二中，所有订单信息可提前一周完整预知，允许在同一自然周内自由组合拼车，且相邻周之间可在满足时效约束的前提下进行跨周合并。
    \item \label{as:single} 问题二仅考虑从南通单一仓库出发的运输场景；问题三中，各仓库独立运作，仓库之间不存在库存调拨或跨仓库订单协同。
    \item \label{as:capacity} 每辆冷藏车的承载能力以托盘数为单位：4.2M型5托，7.6M型11托，9.6M型14托。一条拼车路线内各城市货物托盘数之和不超过所选车型的容量上限。
    \item \label{as:cold} 冷链药品（非"普通药品"类）必须由冷藏车运输；普通药品既可使用常温车，也可使用冷藏车。违反冷链规定的配送被视为不合规。
    \item \label{as:driver} 驾驶员每日驾驶时间不超过10小时，连续行驶允许中途休息，但路线总运输天数由累计驾驶时间和装卸次数共同决定。
    \item \label{as:depot} 所有运输趟次从仓库出发，按规划路线访问各目的地后返回原仓库（闭合路线），不考虑中途补货或临时更改路线。
\end{enumerate}


\subsection{符号约定}

模型中所用符号的含义以及单位见\cref{tab:signal}。
\begin{table}[h]
    \caption{符号约定}\label{tab:signal} \centering
    \begin{tabular}{ccc}
        \toprule[1.5pt]
        符号 & 含义 & 单位 \\
        \midrule[1pt]
        $(\varphi, \lambda)$ & 仓库或城市的经纬度坐标 & $({}^\circ,{}^\circ)$ \\
        $d_{\mathrm{s}}(A,B)$ & 城市 $A$ 与 $B$ 之间的球面直线距离 & km \\
        $R$ & 地球平均半径 ($R = 6371$) & km \\
        $\beta_r$ & 区域 $r$ 的道路绕行系数 & — \\
        $d_{\mathrm{r}}(A,B)$ & 城市 $A$ 与 $B$ 之间的道路估算距离 & km \\
        $\theta_i$ & 从出发仓库看向城市 $i$ 的方位角 & ${}^\circ$ \\
        $\alpha$ & 回程空载油耗系数 ($\alpha = 0.75$) & — \\
        $C_{\mathrm{var}}$ & 单位距离可变成本（油费 + 路桥费）& 元/km \\
        $C_{\mathrm{fix}}$ & 每日固定成本（人工 + 保险 + 审验 + 违章 + 胎耗 + 维保 + 折旧）& 元/天 \\
        $n_v$ & 运输所需车辆数 & 辆 \\
        $T_{ij}$ & 城市 $i$ 到城市 $j$ 的运输天数 & 天 \\
        $p_i$ & 城市 $i$ 的货物托盘数 & 托 \\
        $Q_v$ & 车型 $v$ 的最大承载托数 & 托 \\
        $s_{ij}$ & 合并城市 $i$ 与 $j$ 的CW节约值 & 元 \\
        $\theta_d$ & 方向偏离角阈值 & ${}^\circ$ \\
        $d_{\max}$ & 城市间距离合并上限 & km \\
        $\theta_g$ & 合并组内城市的方位角跨度上限 & ${}^\circ$ \\
        $m_k$ & 路线 $k$ 中访问的城市数 & 个 \\
        $D_k$ & 路线 $k$ 的往返总道路距离 & km \\
        $T_k$ & 路线 $k$ 的运输天数 & 天 \\
        $S_{kj}$ & 仓库 $k$ 在指标 $j$ 上的绝对得分 & 分 \\
        $w_j$ & 指标 $j$ 的熵权权重 & — \\
        $e_j$ & 指标 $j$ 的信息熵 & — \\
        $F_k$ & 仓库 $k$ 的综合得分 & 分 \\
        $F_{\mathrm{sys}}$ & 系统加权综合得分 & 分 \\
        \bottomrule[1.5pt]
    \end{tabular}
\end{table}
